% Part: sets-functions-relations
% Chapter: relations
% Section: relations-as-sets

\documentclass[../../../include/open-logic-section]{subfiles}

\begin{document}

\olfileid{sfr}{rel}{set}
\olsection{Relaties als verzamelingen}

\begin{explain}
In \olref[sfr][set][imp]{sec} hebben we enkele belangrijke
verzamelingen genoemd: $\Nat$, $\Int$, $\Rat$, $\Real$. U herinnert
zich ongetwijfeld nog enkele interessante relaties tussen de !!{element}s
van sommige van deze verzamelingen. Zo heeft elk ervan een gebruikelijke
\emph{ordening}. Daarnaast is er de \emph{identiteit}, die ieder object
met zichzelf en met geen enkel ander object verbindt. We zullen nog veel meer
interessante relaties tegenkomen, en er zijn nog meer relaties mogelijk.
Voordat we die bespreken, merken we echter eerst op dat een relatie als een
bijzonder soort verzameling kan worden opgevat.
  
Hiervoor hebben we twee begrippen uit \olref[sfr][set][pai]{sec} nodig.
Ten eerste is er het \emph{geordende paar}: uit $a$ en $b$ kunnen
we~$\tuple{a, b}$ vormen. De volgorde van de elementen is hier
\emph{wel} van belang. Als $a \neq b$, dan is dus $\tuple{a, b} \neq
\tuple{b, a}$. (Bij ongeordende paren, dat wil zeggen verzamelingen met
$2$ elementen, geldt daarentegen $\{a, b\}=\{b, a\}$.) Ten tweede is er het
\emph{cartesisch product}: als $A$ en $B$ verzamelingen zijn, kunnen we
$A \times B$ vormen, de verzameling van alle paren $\tuple{x, y}$ met
$x \in A$ en $y \in B$. In het bijzonder is $A^{2}= A \times A$ de
verzameling van alle geordende paren uit~$A$.
  
We beschouwen nu één bepaalde relatie op een verzameling: de
$<$-relatie op de verzameling~$\Nat$ van natuurlijke getallen. Neem de
verzameling van alle getallenparen $\tuple{n, m}$ waarvoor $n<m$, dus
\[
  R=\Setabs{\tuple{n, m}}{n, m \in \Nat \text{ en } n<m}.
\]
Dat $n$ kleiner is dan $m$, komt precies overeen met het feit dat het
paar $\tuple{n, m}$ tot $R$ behoort:
\[
      n<m\text{ desda }\tuple{n, m} \in R.
\]
We kunnen de verzameling $R$ dan ook zonder informatieverlies
\emph{opvatten als} de $<$-relatie op $\Nat$.

Op dezelfde manier kunnen we voor iedere relatie tussen getallen een
deelverzameling van $\Nat^{2}$ construeren. Omgekeerd bepaalt iedere
verzameling getallenparen $S \subseteq \Nat^{2}$ een bijbehorende
relatie tussen getallen, namelijk de relatie waarin $n$ tot $m$ staat
desda $\tuple{n, m} \in S$. Dit rechtvaardigt de volgende definitie:
\end{explain}
  
\begin{defn}[Binaire relatie]
Een \emph{binaire relatie} op een verzameling $A$ is een
deelverzameling van~$A^{2}$. Als $R \subseteq A^{2}$ een binaire
relatie op~$A$ is en $x, y \in A$, schrijven we soms $Rxy$ (of $xRy$)
voor $\tuple{x, y} \in R$.
\end{defn}
  
\begin{ex}
  \ollabel{relations}
De verzameling $\Nat^{2}$ van paren natuurlijke getallen kan als volgt
in een tweedimensionale matrix worden weergegeven:
\[
  \begin{array}{ccccc}
  \mathbf{\tuple{ 0,0 }} & \tuple{ 0,1 } &
    \tuple{ 0,2 } & \tuple{ 0,3 } & \ldots\\
  \tuple{ 1,0 } & \mathbf{\tuple{ 1,1 }} &
    \tuple{ 1,2 } & \tuple{ 1,3 } & \ldots\\
  \tuple{ 2,0 } & \tuple{ 2,1 } &
    \mathbf{\tuple{ 2,2 }} & \tuple{ 2,3 } & \ldots\\
  \tuple{ 3,0 } & \tuple{ 3,1 } & \tuple{ 3,2 } &
    \mathbf{\tuple{ 3,3 }} & \ldots\\
  \vdots & \vdots & \vdots & \vdots & \mathbf{\ddots}
  \end{array}
\]
De diagonaal is hier vet gedrukt, want de deelverzameling van $\Nat^2$
die bestaat uit de paren op de diagonaal, namelijk
\[
  \{\tuple{0,0 }, \tuple{ 1,1 }, \tuple{ 2,2 }, \dots\},
  \]
is de \emph{identiteit op}~$\Nat$. (Omdat de identiteit vaak wordt
gebruikt, definiëren we $\Id{A}=\Setabs{\tuple{ x,x }}{x \in A}$ voor
iedere verzameling $A$.) De deelverzameling van alle paren boven de
diagonaal, namelijk
\[
  L = \{\tuple{ 0,1 },\tuple{ 0,2 },\ldots,\tuple{ 1,2 },
  \tuple{ 1,3 }, \dots, \tuple{ 2,3 }, \tuple{ 2,4 },\ldots\},
\]
is de relatie \emph{kleiner dan}, dus $Lnm$ desda $n<m$. De
deelverzameling van alle paren onder de diagonaal, namelijk
\[
  G=\{ \tuple{ 1,0 },\tuple{ 2,0 },\tuple{
    2,1 }, \tuple{ 3,0 },\tuple{ 3,1 },\tuple{ 3,2 }, \dots\},
\]
is de relatie \emph{groter dan}, dus $Gnm$ desda $n>m$. De vereniging
van $L$ met de zojuist gedefinieerde identiteit~$I$, die we $K=L\cup I$
kunnen noemen, is de relatie
\emph{kleiner dan of gelijk aan}: $Knm$ desda $n \le m$. Evenzo is
$H=G \cup I$ de relatie \emph{groter dan of gelijk aan.} De relaties
$L$, $G$, $K$ en $H$ zijn bijzondere relaties die \emph{ordeningen}
heten. $L$ en $G$ hebben de eigenschap dat noch $L$, noch $G$ een getal
met zichzelf in verband brengt (voor alle $n$ geldt dus noch $Lnn$,
noch $Gnn$). Relaties met deze eigenschap heten \emph{irreflexief};
zijn ze bovendien ordeningen, dan heten ze \emph{strikte ordeningen.}
\end{ex}

\begin{explain}
Hoewel ordeningen en de identiteit belangrijke, natuurlijke relaties
zijn, benadrukken we dat volgens onze definitie \emph{iedere}
deelverzameling van $A^{2}$ een relatie op~$A$ is, hoe onnatuurlijk of
gekunsteld die ook lijkt. Zo is $\emptyset$ een relatie op iedere
verzameling (de \emph{lege relatie}, die voor geen enkel elementenpaar
geldt), en is ook $A^{2}$~zelf een relatie op~$A$ (een relatie die voor
ieder paar geldt), de \emph{universele relatie}. Ook iets als
$E=\Setabs{\tuple{n, m}}{n>5 \text{ of } m \times n \ge 34}$ geldt als
een relatie.
\end{explain}
  
\begin{prob}
  Geef de !!{element}s van de relatie $\subseteq$ op de verzameling
  $\Pow{\{a, b, c\}}$.
\end{prob}
\end{document}
